\documentclass[11pt]{article}

\usepackage[a4paper,margin=1in]{geometry}
\usepackage{amsmath,amssymb,amsthm,mathtools,bm}
\usepackage{booktabs}
\usepackage{microtype}
\usepackage{hyperref}
\usepackage[numbers,sort&compress]{natbib}

\newcommand{\Tr}{\operatorname{Tr}}
\newcommand{\dd}{\mathrm{d}}
\newcommand{\ee}{\mathrm{e}}
\newcommand{\EB}{\mathrm{EB}}
\newcommand{\NPT}{\mathrm{NPT}}
\newcommand{\Ccal}{\mathcal C}
\newcommand{\Tcal}{\mathcal T}
\newcommand{\Ocal}{\mathcal O}
\newtheorem{lemma}{Lemma}
\newtheorem{proposition}{Proposition}

\title{A Source-Resolved Quantum Link Budget for Propagating Linearized Gravity}
\author{Anonymous}
\date{August 2026}

\begin{document}
\maketitle

\begin{abstract}
Quantum gravitational radiation, propagating-graviton entanglement, graviton transduction, and gravitational quantum communication have each been studied in increasingly explicit models. What is less often kept visible in one calculation is the complete sequence from a locally initialized quantum source to a remotely accessible quantum record. We construct such a benchmark in weak linearized gravity. A degenerate reference qubit and a branch-common work mode locally prepare opposite coherent amplitudes of a conserved finite-support elastic quadrupole. The two branches have equal total Poincar\'e charges to the working order, so they may be assigned a common first-order asymptotic gravitational dressing; the branch-dependent signal is then the retarded field generated by their conserved difference stress tensor. At an exact controller-empty handoff, all branch dependence in the linear coherent network can be compressed into a single virtual bosonic difference mode. For receiver states evaluated after this handoff, including the coherent effect of any precursor radiation emitted during the encoder, the source-to-memory transfer factorizes as
\begin{equation}
\tau_c(t)=
\underbrace{\frac{\kappa_{g,A}}{\kappa_A}}_{\text{source gravitational branching}}
\underbrace{\frac{25\Ocal}{16(kR)^2}}_{\text{free-space mode capture}}
\underbrace{\frac{\kappa_{g,B}}{\kappa_B}}_{\text{receiver gravitational branching}}
\underbrace{\Tcal_f(t)}_{\text{temporal loading}},
\qquad 0\le \Tcal_f\le1.
\label{eq:abstract-link}
\end{equation}
Receiver and source Gaussian noise define a memory quantum excess $\Delta_{\rm mem}=\tau_c-m_c$, while a separate readout channel gives the accessible excess $\Delta_{\rm acc}=\tau_r\Delta_{\rm mem}-m_r$. Established bosonic-channel criteria imply non-entanglement-breaking transmission when the corresponding excess is positive. For vacuum pure loss we also obtain the exact reference--receiver negativity and show that its optimum obeys $\mathcal N_{\max}=\tau_c-2\tau_c^{3/2}+O(\tau_c^2)$ in the weak-link limit. A deliberately aggressive kilogram--meter--MHz benchmark gives a coherent link ceiling of order $10^{-42}$ when both matter--gravity interfaces are ordinary passive mechanical modes, compared with order $10^{-22}$ if one interface is idealized. The result is a source-resolved normalization and capability benchmark, not a near-term experimental proposal and not a new Gaussian-channel theorem.
\end{abstract}

\section{Introduction}

Several distinct questions are often grouped under the heading of quantum communication through gravity. One may ask how quantum matter radiates into a quantized gravitational field, how propagating gravitons correlate separated matter systems, how a quantum gravitational-wave mode is absorbed by a detector, or whether an effective gravitational interaction channel can preserve entanglement. Each ingredient now has substantial prior literature.

Closed or conserved branch-dependent sources and their emitted coherent graviton states have been analyzed directly, including the relation between the coherent-state distance of two radiation histories and vacuum which-path decoherence \cite{Matsui2026}. Classical conserved stress tensors driving the quantized gravitational field produce coherent states whose expectation value reproduces the retarded classical waveform \cite{LagaSuyama2026}. Propagating gravitons have been used to generate delayed matter entanglement \cite{TrengganaZen2026}, while resonant matter systems have been treated as receivers of nonclassical gravitational-wave states with open-system noise included \cite{ToccaceloEtAl2026,TobarEtAl2024}. Gravitational quantum communication has also been phrased directly in terms of entanglement-breaking channels and thermal thresholds \cite{MariEtAl2026,ToccaceloAndersenBrask2025,MikiLiChen2026}. Quantum photon--graviton conversion, cavity-QED transduction, and stimulated graviton absorption/emission have likewise been proposed \cite{IkedaEtAl2025,AraniLamineSoda2026,Schutzhold2025}.

Our question is narrower. We ask what happens when the interfaces usually idealized away are kept explicit in one source-to-readout normalization:
\begin{equation}
\boxed{
\text{local equal-charge encoding}
\to
\text{conserved radiator}
\to
\text{gravitational output port}
\to
\text{retarded free-space mode}
\to
\text{receiver memory}
\to
\text{readout}.}
\label{eq:chain}
\end{equation}
The objective is not to introduce a new information-theory criterion. The entanglement-breaking structure of one-mode Gaussian channels is established \cite{Holevo2008EB}, as is the sensitivity of non-Gaussian entanglement to attenuation and amplification \cite{FilippovZiman2014}. The contribution sought here is instead a source-resolved gravitational link calculation in which the physical meaning and normalization of each efficiency and noise term remain visible.

A second issue is locality. Gauge-invariant matter operators in gravity require nonlocal gravitational dressing, so the ordinary local-QFT picture of exactly commuting compact subsystems cannot simply be imported into this problem \cite{DonnellyGiddings2016}. Perturbative gravitational splitting nevertheless permits information inside a region to be hidden from exterior measurements, at first order, except through total Poincar\'e charges \cite{DonnellyGiddings2018}. We therefore encode the source qubit into two branches with equal total Poincar\'e charges and use a common asymptotic dressing. The operational signal is the later retarded difference field, not an assumed instantaneous branch-dependent Coulomb dressing.

The resulting paper has five logical layers. First, we construct the equal-charge conserved source. Second, we identify the time at which its locally generated bosonic code becomes a fixed channel input. Third, we derive a four-factor coherent link budget. Fourth, we include source and receiver noise and quantify actual entanglement transfer. Fifth, we append a separate noisy readout stage so that memory capture is not confused with accessible quantum reception.

\section{Regime and explicit conserved source}
\label{sec:source}

We work in linearized gravity around Minkowski spacetime, with compact nonrelativistic elastic source and receiver modes. For a mechanical device of scale $L$, resonance frequency $\omega$, and longitudinal sound speed $c_s$, define
\begin{equation}
q\equiv\frac{\omega L}{c_s},
\qquad
\beta\equiv\frac{\omega L}{c},
\qquad
\Ccal\equiv\frac{2GM}{c^2L}.
\end{equation}
The controlled mechanical regime is
\begin{equation}
\frac{|u|}{L}\ll1,
\qquad q\ll1,
\qquad \beta\ll1,
\qquad \Ccal\ll1.
\label{eq:regime}
\end{equation}
We additionally use a narrowband rotating-wave/input-output model around one source and one receiver quadrupole resonance, with $g,\kappa_A,\kappa_B,1/T\ll\omega$ whenever the corresponding rates are used.

\subsection{Four-spoke plus mode}

The explicit source consists of four endpoint masses $\mu$ connected to a central compact hub by four identical longitudinal elastic spokes of reference length $L$. The spokes lie along the $\pm x$ and $\pm y$ axes. The plus mode expands the $x$ pair while contracting the $y$ pair, with the opposite branch obtained by reversing the sign of the displacement coordinate $u$.

For a spoke mode normalized to unit endpoint displacement,
\begin{equation}
f_q(x)=\frac{\sin(qx/L)}{\sin q},
\end{equation}
and the endpoint traction condition fixes the spoke mass $m_r$ as
\begin{equation}
\boxed{\frac{m_r}{\mu}=q\tan q.}
\label{eq:mrmu}
\end{equation}
Including endpoint and spoke rest mass, the branch-difference STF quadrupole is
\begin{equation}
\boxed{
\Delta Q_{xx}=8\mu Lu\frac{\tan q}{q},
\qquad
\Delta Q_{yy}=-\Delta Q_{xx},
\qquad
\Delta Q_{zz}=0.}
\label{eq:DeltaQ}
\end{equation}
The generalized mode mass is
\begin{equation}
\boxed{
M_{\rm eff}=4\mu A(q),
\qquad
A(q)=\frac12+\frac{q}{\sin2q}.}
\label{eq:Meff}
\end{equation}
Quantizing $u=u_{\rm zpf}(a+a^\dagger)$ gives
\begin{equation}
q_{01}=4\mu L\frac{\tan q}{q}
\sqrt{\frac{\hbar}{2M_{\rm eff}\omega}},
\end{equation}
and the intrinsic plus-mode spontaneous graviton linewidth
\begin{equation}
\boxed{
\kappa_g(q)=
\frac{8G\mu L^2\omega^4}{5c^5}\,
\Ccal_\kappa(q),
\qquad
\Ccal_\kappa(q)=\frac{(\tan q/q)^2}{A(q)}.}
\label{eq:kappag}
\end{equation}
For $q\ll1$,
\begin{equation}
\Ccal_\kappa(q)=1+\frac{q^2}{3}+O(q^4).
\end{equation}
The full source stress tensor includes endpoints, spokes, hub, and controller and is conserved to the working elastic/weak-field order. Internal stress is therefore part of the same conserved source; it is not an independent radiative contribution to be appended to the mass quadrupole after the fact.

For the endpoint coordinates $X=L+u$ and $Y=L-u$, the branch-carrying plus combination obeys
\begin{equation}
Q_{xx}-Q_{yy}=2\mu(X^2-Y^2)=8\mu Lu
\end{equation}
exactly. The leading $u^2$ geometric term is branch even and occupies the dc/$2\omega$ sector rather than the branch-odd resonant mode.

\section{Local equal-charge encoding and gravitational dressing}
\label{sec:encoding}

Let $S$ be a degenerate internal reference qubit, $a$ the source plus mode, and $w$ a branch-common work mode. We use the local resonant encoder
\begin{equation}
\boxed{
H_{\rm enc}=\hbar g\sigma_z(a^\dagger w+a w^\dagger).}
\label{eq:Henc}
\end{equation}
The initial state is taken as
\begin{equation}
|\Psi_0\rangle
=\frac{|0\rangle_S+|1\rangle_S}{\sqrt2}
\otimes|0\rangle_a\otimes|\zeta\rangle_w\otimes|0\rangle_E,
\label{eq:init}
\end{equation}
where $E$ denotes the source output ports, including the graviton vacuum.

In the lossless limit, branch $s=\pm1$ evolves as
\begin{equation}
\alpha_s(t)=-is\zeta\sin(gt),
\qquad
\gamma_w(t)=\zeta\cos(gt).
\end{equation}
With source-mode amplitude damping rate $\kappa_A$ define
\begin{equation}
\Omega=\sqrt{g^2-\kappa_A^2/16}.
\end{equation}
Then
\begin{align}
\gamma_w(t)&=\zeta\ee^{-\kappa_A t/4}
\left[\cos(\Omega t)+\frac{\kappa_A}{4\Omega}\sin(\Omega t)\right],\\
\alpha_s(t)&=-is\zeta\frac{g}{\Omega}\ee^{-\kappa_A t/4}\sin(\Omega t).
\end{align}
The first exact controller-empty time is
\begin{equation}
\boxed{
T_*=\frac{\pi-\arctan(4\Omega/\kappa_A)}{\Omega},}
\label{eq:Tstar}
\end{equation}
for which
\begin{equation}
\gamma_w(T_*)=0,
\qquad
\alpha_s(T_*)=-is\zeta\ee^{-\kappa_A T_*/4}.
\end{equation}
For coherent inputs and vacuum linear loss ports, the conditional network remains a product of coherent states, so the work mode is exactly branch common and factorized at the handoff.

\subsection{Equal-charge code}

The mirrored plus branches are chosen to have, to the working order,
\begin{equation}
\boxed{\Delta P^\mu=0,\qquad \Delta M^{\mu\nu}=0.}
\label{eq:equalcharges}
\end{equation}
The total energy is branch even, inversion symmetry cancels net momentum, the radial plus motion carries no net angular momentum, and the center of energy remains at the hub. The reference qubit is degenerate and the ideal work-mode energy and momentum are branch common.

Equation~\eqref{eq:equalcharges} matters because physical matter operators in gravity require gravitational dressing \cite{DonnellyGiddings2016}. At first perturbative order, the gravitational-splitting construction allows states with the same total Poincar\'e charges to share a standard asymptotic dressing whose exterior observables are insensitive to the encoded internal information \cite{DonnellyGiddings2018}. We therefore use a common asymptotic dressing for the two code states. Their quadrupole difference is not an asymptotic Poincar\'e charge; its later physical effect is carried by the retarded field sourced by the conserved difference stress tensor $\Delta T^{\mu\nu}$.

This statement is deliberately perturbative and code restricted. We do not claim exact gravitational tensor-factor locality, nor do we claim that states with different total mass or momentum are exterior-indistinguishable.

\subsection{Causal origin versus channel handoff}

There are two different clocks. The local causal intervention begins at $t=t_s$. The fixed bosonic input code is established only at $t=t_s+T_*$. During the encoder, the bosonic branch separation is still being created. In the lossless limit,
\begin{equation}
N_{\Delta,{\rm all}}(t)=4|\zeta|^2\sin^2(gt).
\end{equation}
At the controller-empty handoff, including precursor radiation emitted during the encoder,
\begin{equation}
\boxed{
N_{\Delta,A}(T_*)+
\sum_jN_{\Delta,j}^{\rm enc}=4|\zeta|^2.}
\label{eq:distance-conservation}
\end{equation}
After this time the total branch-distance norm is fixed and branch-independent linear dynamics only redistribute it among modes. The causal clock therefore starts at $t_s$, but the fixed-input bosonic channel used below begins at $T_*$.

The precursor is not discarded. The normalized gravitational waveform used later contains both radiation emitted during the encoder and radiation emitted after handoff. If the source--receiver separation is small enough, some precursor can reach and begin driving the receiver before $T_*$. During that early interval the physical process is a controlled qubit-to-multimode state-generation problem, not yet a fixed-input one-mode channel; its entanglement should be evaluated directly from the instantaneous coherent branch state. For any receiver state evaluated after handoff, however, the fixed global virtual mode defined at $T_*$ includes all precursor modes already in flight, and the receiver amplitude includes the coherent effect of any precursor that has already arrived.

For a centrally triggered finite source, a source point $\mathbf x$ can acquire branch dependence only after a causal internal signal reaches it. A later signal reaching receiver point $\mathbf y$ satisfies
\begin{equation}
t_{\rm arr}-t_s\ge
\frac{|\mathbf x-\mathbf x_0|+|\mathbf y-\mathbf x|}{c}
\ge\frac{|\mathbf y-\mathbf x_0|}{c},
\end{equation}
so the center-origin $R/c$ lower bound remains valid. Before the receiver's causal past intersects the branch-dependent intervention history, physical receiver observables act proportionally to the identity on the equal-charge code to the working perturbative order. Equivalently, the encoded source-to-receiver map is a replacer channel on that code before causal arrival.

\section{Virtual difference mode after handoff}
\label{sec:virtualmode}

The local encoder does not assume a bosonic mode already entangled with $S$. That bosonic channel input emerges only after the local preparation has completed.

\begin{lemma}[Virtual difference-mode reduction]
After the handoff, suppose the branch-conditioned state of all bosonic degrees of freedom in the vacuum/coherent linear network is, after removing branch-common displacements,
\begin{equation}
|\Psi_s\rangle=|s\bm{\alpha}\rangle,
\qquad s=\pm1.
\end{equation}
Let $A=\|\bm{\alpha}\|$ and define
\begin{equation}
\boxed{
d=\frac{1}{A}\sum_j\alpha_j^*b_j.}
\label{eq:virtualmode}
\end{equation}
Then a branch-independent passive mode rotation maps
\begin{equation}
|s\bm{\alpha}\rangle
\longmapsto
|sA\rangle_d\otimes|0\rangle_\perp.
\label{eq:modecompression}
\end{equation}
For any selected receiver memory mode with coherent branch amplitude $\alpha_B$, the reduced map from $d$ to that memory is a pure-loss projection of transmissivity
\begin{equation}
\boxed{
\eta_B=\frac{|\alpha_B|^2}{A^2}
=\frac{N_{\Delta,B}}{N_{\Delta,{\rm all}}}.}
\label{eq:etaB}
\end{equation}
\end{lemma}

The proof is simply a unitary rotation in one-particle mode space that aligns the displacement vector $\bm{\alpha}$ with one basis vector. Because coherent states transform covariantly under passive mode rotations, all branch dependence is compressed into $d$ and the orthogonal modes are vacuum. This lemma is the formal bridge between the physical local encoder and the standard binary-coherent bosonic-channel problem.

A branch-common controller or background gravitational coherent field causes only a common displacement, $|\gamma_C\pm\alpha\rangle$, and is removed by the same branch-independent displacement. It does not change $N_\Delta$, NPT, or entanglement-breaking status. Controller fluctuations above the coherent-state vacuum level are different: they are genuine added noise and must be propagated through the channel.

\section{Source gravitational branching and free-space propagation}
\label{sec:sourceprop}

Let the source mode couple to independent narrowband Markov output ports $j$ with rates $\kappa_j$. For a mirrored coherent trajectory $\alpha_s(t)=s\alpha(t)$, the branch-difference output amplitude is
\begin{equation}
\Delta b_j^{\rm out}(t)=2\sqrt{\kappa_j}\,\alpha(t).
\end{equation}
Hence
\begin{equation}
N_{\Delta,j}=4\kappa_j\int\dd t\,|\alpha(t)|^2.
\end{equation}
The fraction of the complete branch-distance norm emitted into the gravitational port is therefore
\begin{equation}
\boxed{
\beta_{g,A}
\equiv\frac{N_{\Delta,g}}{N_{\Delta,{\rm tot}}}
=\frac{\kappa_{g,A}}{\kappa_A},}
\qquad
\kappa_A=\sum_j\kappa_j.
\label{eq:betaA}
\end{equation}
This equality is a narrowband amplitude-damping result. If the port couplings vary appreciably across the source waveform bandwidth, the correct object is the waveform-weighted ratio
\begin{equation}
\beta_{g,A}[\alpha]
=
\frac{\int\dd\Omega\,\kappa_g(\omega+\Omega)|\widetilde\alpha(\Omega)|^2}
{\int\dd\Omega\,\kappa_{\rm tot}(\omega+\Omega)|\widetilde\alpha(\Omega)|^2}.
\label{eq:freqbeta}
\end{equation}

Equation~\eqref{eq:betaA} is not a generic coherence parameter. Pure phase diffusion can destroy the source-reference entanglement without changing the source energy or any branching ratio. Non-Gaussian source dephasing must therefore be represented by a separate source channel rather than hidden inside $\beta_{g,A}$.

\subsection{Wave-zone mode capture}

For aligned compact plus quadrupoles in the wave zone, the retarded source-to-receiver normalization gives the useful receiver gravitational-port rate
\begin{equation}
\kappa_\Delta(R)=\eta_{\rm store}(R)\kappa_{g,B},
\end{equation}
with
\begin{equation}
\boxed{
\eta_{\rm store}(R)=\frac{25\Ocal}{16(kR)^2},
\qquad k=\omega/c,}
\label{eq:etastore}
\end{equation}
where $0\le\Ocal\le1$ collects residual temporal, tensor, polarization, and orientation mode overlap.

The coefficient has a useful reciprocal-antenna interpretation. The on-axis directivity of the aligned plus source is $G_A=5/2$, while the critically coupled $l=2$ receiver effective area $A_e=5\pi/(2k^2)$ corresponds to reciprocal gain $G_B=5/2$. Thus
\begin{equation}
\boxed{
\eta_{\rm store}
=\Ocal G_AG_B\left(\frac{\lambda}{4\pi R}\right)^2
=\frac{25\Ocal}{16(kR)^2}.}
\label{eq:friis}
\end{equation}
This is ordinary reciprocal far-field link physics specialized to the quadrupole channel, not a new antenna theorem. Critical-coupling and temporal coupled-mode normalization are standard \cite{RuanFan2012}.

\section{Receiver memory and the four-factor coherent link}
\label{sec:memory}

Let the receiver memory mode $b$ have total linewidth $\kappa_B$ and intrinsic gravitational linewidth $\kappa_{g,B}$. Define
\begin{equation}
\boxed{\beta_{g,B}=\frac{\kappa_{g,B}}{\kappa_B}.}
\label{eq:betaB}
\end{equation}
After retarded arrival, let $f(t)$ be the normalized temporal envelope of the selected incoming source mode,
\begin{equation}
\int_0^\infty\dd t\,|f(t)|^2=1.
\end{equation}
Here $f$ is the complete source gravitational waveform: it may contain encoder precursor radiation emitted before $T_*$ as well as later source radiation. When a receiver state is evaluated at a time after handoff, the convolution below includes any precursor contribution that has already arrived.

The linear receiver equation contains the selected gravitational input with rate $\kappa_\Delta=\eta_{\rm store}\kappa_{g,B}$. The coherent loading from the selected mode is therefore
\begin{equation}
\tau_{f,B}(t)
=\kappa_\Delta
\left|
\int_0^t\dd s\,
\ee^{-\kappa_B(t-s)/2}f(s)
\right|^2.
\end{equation}
Define the dimensionless temporal loading factor
\begin{equation}
\boxed{
\Tcal_f(t)
=\kappa_B
\left|
\int_0^t\dd s\,
\ee^{-\kappa_B(t-s)/2}f(s)
\right|^2,}
\qquad
0\le\Tcal_f(t)\le1.
\label{eq:Tf}
\end{equation}
The upper bound follows from Cauchy--Schwarz and is saturated at a chosen target time by the time reverse of the receiver impulse response.

Combining the post-handoff source branching, propagation, receiver branching, and temporal loading gives the central coherent-transfer result:
\begin{proposition}[Post-handoff gravitational quantum-link budget]
For the vacuum/coherent, narrowband linear amplitude-damping model, let the fixed virtual difference mode be defined at $T_*$. For any receiver state evaluated at $t\ge T_*$, with the complete waveform $f$ including any encoder precursor, the receiver branch-distance fraction relative to the completed code is
\begin{equation}
\boxed{
\tau_c(t)=
\beta_{g,A}\,
\eta_{\rm store}(R)\,
\beta_{g,B}\,
\Tcal_f(t).}
\label{eq:master}
\end{equation}
Defining
\begin{equation}
\boxed{\eta_Q^{\rm link}=\beta_{g,A}\eta_{\rm store}\beta_{g,B},}
\label{eq:etalink}
\end{equation}
one has $\tau_c(t)=\eta_Q^{\rm link}\Tcal_f(t)$ and $\tau_c\le\eta_Q^{\rm link}$.
\end{proposition}

The factors in Eq.~\eqref{eq:master} are not duplicate normalizations. The same product follows from an independent cascaded bosonic network: a source beamsplitter of transmissivity $\beta_{g,A}$, a propagation beamsplitter of transmissivity $\eta_{\rm store}$, and a receiver loading channel whose coherent fraction is $\beta_{g,B}\Tcal_f$. The product is therefore the canonical serial loss of the post-handoff virtual mode.

For $t<T_*$ the same expression, if normalized to the eventual code norm $4|\zeta|^2$, should be interpreted only as a coherent response coefficient. The instantaneous global coherent branch state can still be compressed into an instantaneous virtual mode, but the logical-to-bosonic code itself is still being generated and the fixed-input channel theorem is not yet the appropriate description.

Several waveform results are now merely choices of the order-unity factor $\Tcal_f$. A matched passive exponential gives
\begin{equation}
\Tcal_{\rm exp}^{\max}=4\ee^{-2}\simeq0.541341,
\end{equation}
a previously analyzed smooth $\sin^4$ pulse gives
\begin{equation}
\Tcal_{\sin^4}^{\max}\simeq0.7980213,
\end{equation}
and target-time optimization yields
\begin{equation}
\Tcal_{\rm opt}(t)=1-\ee^{-\kappa_B t}\to1.
\end{equation}
Temporal shaping can remove an order-unity mismatch penalty, but it cannot exceed the source-branching, propagation, or receiver-branching ceiling.

\section{Noise, quantum excess, and accessible readout}
\label{sec:noise}

At a fixed time, the reduced selected-mode-to-memory map is phase-insensitive Gaussian in the linear Markov model. Let $m_A$ denote Gaussian source noise in the virtual source channel convention and let $m_B(t)$ denote receiver occupation generated by the initial receiver state and occupied baths. For an initially cooled receiver with Markov bath injection rate $\Gamma_{\rm th}$,
\begin{equation}
m_B(t)=n_0\ee^{-\kappa_B t}
+\frac{\Gamma_{\rm th}}{\kappa_B}
\left(1-\ee^{-\kappa_B t}\right).
\label{eq:mB}
\end{equation}
The memory-channel parameters are
\begin{align}
\boxed{\tau_c(t)}&=\beta_{g,A}\eta_{\rm store}\beta_{g,B}\Tcal_f(t),\\
\boxed{m_c(t)}&=m_B(t)
+\eta_{\rm store}\beta_{g,B}\Tcal_f(t)m_A.
\label{eq:memoryparams}
\end{align}
Define the memory quantum excess
\begin{equation}
\boxed{\Delta_{\rm mem}(t)=\tau_c(t)-m_c(t).}
\label{eq:Dmem}
\end{equation}
For this phase-insensitive Gaussian family, established channel theory gives
\begin{equation}
\boxed{\Delta_{\rm mem}>0\quad\Longleftrightarrow\quad\text{memory channel is non-EB}.}
\label{eq:memnonEB}
\end{equation}
For the binary-coherent hybrid input produced by the virtual-mode reduction, the same boundary is the NPT boundary. We use this as prior channel theory rather than as a novelty claim.

Equation~\eqref{eq:memoryparams} can be written as
\begin{equation}
\eta_{\rm store}\beta_{g,B}\Tcal_f
\left(\beta_{g,A}-m_A\right)>m_B
\label{eq:memcondition}
\end{equation}
for memory quantum capability. Non-Gaussian source dephasing or nonlinear receiver noise is not represented by the single scalar $m$ and must be analyzed as an additional channel.

\subsection{Accessible readout}

A receiver memory is not yet an accessible register. Let the memory be followed by a phase-insensitive readout channel with parameters $(\tau_r,m_r)$. Channel composition gives
\begin{equation}
\boxed{\tau_{\rm acc}=\tau_r\tau_c,}
\qquad
\boxed{m_{\rm acc}=m_r+\tau_r m_c.}
\label{eq:readoutcompose}
\end{equation}
Define
\begin{equation}
\boxed{\Delta_{\rm acc}=\tau_{\rm acc}-m_{\rm acc}.}
\end{equation}
Then
\begin{equation}
\boxed{
\Delta_{\rm acc}=\tau_r\Delta_{\rm mem}-m_r.}
\label{eq:Dacc}
\end{equation}
Thus accessible quantum capability requires
\begin{equation}
\boxed{\tau_r\Delta_{\rm mem}>m_r.}
\label{eq:accesscond}
\end{equation}
A strong gravitational absorber can therefore be a poor accessible quantum receiver if its readout destroys the memory's small positive quantum excess.

\section{Exact pure-loss entanglement amount}
\label{sec:negativity}

The sign condition alone is insufficient for an extremely weak gravitational link. In vacuum, for a receiver state evaluated after handoff, Eq.~\eqref{eq:master} is a pure-loss projection of the fixed virtual mode with transmissivity
\begin{equation}
\eta\equiv\tau_c.
\end{equation}
Consider the post-handoff virtual-mode state
\begin{equation}
|\Phi_A\rangle
=\frac{|0\rangle_S|+A\rangle_d+|1\rangle_S|-A\rangle_d}{\sqrt2}.
\label{eq:hybrid}
\end{equation}
After pure loss, the receiver branch amplitude is $B=\sqrt\eta A$ and the complementary environment amplitude is $E=\sqrt{1-\eta}A$. Define the coherent-state overlaps
\begin{equation}
s_B=\ee^{-2\eta A^2},
\qquad
s_E=\ee^{-2(1-\eta)A^2}.
\end{equation}
The partial transpose of the exact reference--receiver state has one negative eigenvalue, giving
\begin{equation}
\boxed{
\mathcal N(\eta,A)=
\frac{
\sqrt{(1+s_E)^2-4s_Es_B^2}
-(1-s_E)}{4}.}
\label{eq:exactN}
\end{equation}
For every $A>0$ and $\eta>0$, Eq.~\eqref{eq:exactN} is positive. At fixed finite $A$ and $\eta\ll1$,
\begin{equation}
\boxed{
\mathcal N(\eta,A)=
\eta\frac{2A^2}{\ee^{2A^2}-1}+O(\eta^2).}
\label{eq:fixedAweak}
\end{equation}
Optimizing the branch amplitude as $\eta\to0$ gives
\begin{equation}
\boxed{A_{\rm opt}^2\sim\sqrt\eta,}
\end{equation}
and
\begin{equation}
\boxed{
\mathcal N_{\max}(\eta)
=\eta-2\eta^{3/2}+\frac{13}{3}\eta^2+O(\eta^{5/2}).}
\label{eq:Nopt}
\end{equation}
The physically relevant conclusion is therefore not merely that entanglement survives mathematically for every nonzero vacuum link. The best deliverable negativity is asymptotically linear in the complete end-to-end transmissivity.

\section{Passive-matter ceiling and active loopholes}
\label{sec:bounds}

For a stationary passive nonrelativistic source or receiver described by an ordinary coordinate quadrupole, an energy-weighted quadrupole sum rule bounds the positive absorptive gravitational spectral weight. The derivation is given in Appendix~\ref{app:sumrule}; the result is
\begin{equation}
\boxed{
\frac{\kappa_g}{\omega}
\lesssim\frac23\Ccal\beta^3.}
\label{eq:sumrule}
\end{equation}
If $Q=\omega/\kappa_{\rm tot}$, this implies the gravitational branching bound
\begin{equation}
\boxed{
\beta_g\lesssim
\min\left[1,\frac23Q\Ccal\beta^3\right].}
\label{eq:betabound}
\end{equation}
When both endpoints lie in this passive nonrelativistic class,
\begin{equation}
\boxed{
\eta_Q^{\rm link}
\lesssim
\frac{25\Ocal}{16(kR)^2}
\prod_{j=A,B}
\min\left[1,\frac23Q_j\Ccal_j\beta_j^3\right].}
\label{eq:passivebound}
\end{equation}
In the unsaturated regime,
\begin{equation}
\boxed{
\eta_Q^{\rm link}
\lesssim
\frac{25\Ocal}{36(kR)^2}
Q_AQ_B\Ccal_A\Ccal_B\beta_A^3\beta_B^3.}
\label{eq:passiveunsat}
\end{equation}
This is explicitly a passive, nonrelativistic matter bound. It is not asserted for active/inverted many-body states, relativistic QFT receivers, or strongly self-gravitating systems.

Known correlated matter states can exhibit collective gravitational transition rates with $N^2$ scaling \cite{QuinonesEtAl2017}. Such enhancement can be useful if ordinary internal loss is the dominant bottleneck. If a factor $F$ multiplies all gravitational receiver rates while a nongravitational internal loss rate $\kappa_i$ remains fixed, then
\begin{equation}
\boxed{
\beta_{\rm useful}(F)
=\beta_{\rm mode}
\frac{F\kappa_{g,0}}{F\kappa_{g,0}+\kappa_i}
\longrightarrow\beta_{\rm mode}.}
\label{eq:activebeta}
\end{equation}
Collectivity can therefore remove internal-loss suppression and shorten the interaction time, but if all gravitational modes are enhanced equally it saturates at the same source-mode selectivity. Moreover, the favorable collective states also have enhanced gravitational vacuum transition dynamics. A complete active receiver must therefore include its pump, spontaneous noise, and non-Gaussian state preparation rather than replacing $\kappa_g$ by $N^2\kappa_g$ in Eq.~\eqref{eq:master}.

\section{Benchmark and hierarchy of bottlenecks}
\label{sec:benchmark}

To expose the source/receiver distinction, consider the deliberately aggressive benchmark
\begin{equation}
M_e=4\,{\rm kg},\quad
L=1\,{\rm m},\quad
f=1\,{\rm MHz},\quad
Q=10^{12},\quad
kR=10,\quad
\Ocal=1.
\end{equation}
For the explicit four-spoke leading-order mode,
\begin{equation}
\beta_g\simeq1.09386\times10^{-20},
\qquad
\eta_{\rm store}=1.5625\times10^{-2}.
\end{equation}
The coherent link ceiling is therefore
\begin{table}[h]
\centering
\begin{tabular}{lcc}
\toprule
Interfaces & $(\beta_{g,A},\beta_{g,B})$ & $\eta_Q^{\rm link}$ \\
\midrule
ordinary source + ordinary receiver & $(\beta_g,\beta_g)$ & $1.87\times10^{-42}$ \\
one ideal gravitational interface & $(1,\beta_g)$ or $(\beta_g,1)$ & $1.71\times10^{-22}$ \\
both ideal gravitational interfaces & $(1,1)$ & $1.5625\times10^{-2}$ \\
\bottomrule
\end{tabular}
\caption{Separation of source and receiver gravitational interface penalties at fixed $kR=10$ and perfect residual mode overlap. ``Ideal'' means unit gravitational branching, not a claimed realizable device.}
\label{tab:benchmark}
\end{table}
A matched passive exponential multiplies the ordinary/ordinary ceiling by $4\ee^{-2}$, giving
\begin{equation}
\tau_{\max}\simeq1.01\times10^{-42}.
\end{equation}
Equation~\eqref{eq:Nopt} then shows that the maximum pure-loss negativity is of the same order. The receiver-local order-$10^{-22}$ scale corresponds to starting with an already normalized incoming graviton wavepacket; the true ordinary-matter source contributes a second gravitational branching penalty.

The underlying intrinsic gravitational linewidth is approximately
\begin{equation}
\kappa_g\simeq6.87\times10^{-26}\,{\rm s}^{-1},
\end{equation}
so a purely gravitational passive source would have a characteristic lifetime
\begin{equation}
\kappa_g^{-1}\simeq4.6\times10^{17}\,{\rm yr}.
\end{equation}
This illustrates the speed--efficiency tension. Adding ordinary dissipative source broadening shortens the pulse but reduces the gravitational branching fraction. Coherent shaping can improve temporal matching without adding that extra lossy port, but it cannot change the branch-distance fraction associated with fixed physical output couplings.

\section{Controlled corrections and limitations}
\label{sec:limitations}

The link budget is deliberately a leading-order benchmark. The following qualifications are important.

\paragraph{Finite source size.}
The first finite-wavelength correction to the explicit quadrupole linewidth has the form
\begin{equation}
\kappa_g(q,\beta)=\kappa_g^{(Q)}(q)
\left[1-\frac{2a(q)}7\beta^2+O(\beta^4)\right],
\end{equation}
where
\begin{equation}
a(q)=\frac12+\frac{\cot q}{q}-\frac1{q^2}.
\end{equation}
For $q\to0$, this becomes $\kappa_g=\kappa_g^{(Q)}[1-\beta^2/21+O(\beta^4)]$.

\paragraph{Controller radiation.}
Branch-common coherent controller radiation is a common displacement and does not alter the branch-distance norm or entanglement. Thermal or nonclassical controller fluctuations are genuine channel noise.

\paragraph{Source dephasing.}
The scalar $\beta_{g,A}$ tracks coherent amplitude-damping branching, not arbitrary phase coherence. Complete phase randomization can make the reference--source state separable even with unit gravitational energy branching. Non-Gaussian dephasing must be included as a separate source map.

\paragraph{Frequency dependence.}
Constant branching fractions require a narrowband Markov regime. Outside it, Eq.~\eqref{eq:freqbeta} or the corresponding full spectral transfer kernel must replace the scalar ratios.

\paragraph{Reciprocal feedback.}
For a weak one-way link, the loop correction is small. In the audited wave-zone model,
\begin{equation}
|L(\nu)|\le4\eta_{\rm store}\beta_{g,A}\beta_{g,B},
\end{equation}
and the first source-controlled round-trip echo appears only after approximately $3R/c$. The one-way channel is therefore the leading term in a controlled weak-feedback expansion.

\paragraph{Gravity-specific locality.}
The pre-arrival replacer statement is restricted to the equal-Poincar\'e-charge encoded subspace and to the working perturbative gravitational-splitting framework. It is not an assertion of exact local tensor-factor structure in quantum gravity.

\paragraph{Readout.}
Equation~\eqref{eq:master} terminates at the receiver memory. Operational communication or measurement requires the additional readout condition in Eq.~\eqref{eq:accesscond}.

\section{Discussion}

The value of Eq.~\eqref{eq:master} is not that products of efficiencies are unfamiliar. Its value is that, for this gravitational problem, the usually separated physical descriptions can be joined without silently normalizing away one of the hard interfaces.

The source interface asks what fraction of an explicitly locally prepared branch record enters the gravitational output rather than uncontrolled mechanical ports. The propagation factor asks how much of that normalized gravitational mode is geometrically and tensorially accessible to the remote receiver. The receiver interface asks what fraction of its damping belongs to gravity rather than ordinary material loss. Temporal loading asks whether the actual fixed source waveform matches the receiver. Noise then determines whether the memory channel is genuinely non-entanglement-breaking, and readout determines whether that small quantum excess survives into an accessible register.

This hierarchy changes how proposed improvements should be judged. Pulse shaping can improve $\Tcal_f$ only up to unity. Directional engineering can improve mode overlap only up to the normalized capture ceiling. Collective or active matter can improve gravitational branching relative to fixed internal loss, but it must also account for the associated gravitational vacuum dynamics and pump noise. A proposal that increases a classical response amplitude without increasing the complete channel's distance from the entanglement-breaking boundary has not necessarily improved quantum information transfer.

The equal-charge code is also essential conceptually. Because the two source branches share their total Poincar\'e charges, the branch label is not encoded in the asymptotic Coulomb data that gravitational dressing must carry. The branch-sensitive information resides in internal multipole structure and is exposed to the remote receiver through retarded dynamics. This is the appropriate perturbative replacement for an overly strong claim of compact local subsystems in gravity.

Finally, the exact pure-loss result makes the feasibility interpretation quantitative. In a weak link, the maximum reference--receiver negativity scales as the link transmissivity itself. Thus a fantastically small link is not rescued by choosing an arbitrarily large coherent branch separation; the optimum instead moves toward a weak branch amplitude and the deliverable entanglement remains linear in the full serial efficiency.

\section{Conclusion}

We have constructed a weak-field source-to-readout gravitational quantum link in which the local source, conserved support, gravitational emission, retarded propagation, resonant memory, noise, and readout are all explicitly separated. The central post-handoff coherent-transfer coefficient is
\begin{equation}
\boxed{
\tau_c(t)=
\beta_{g,A}\eta_{\rm store}(R)\beta_{g,B}\Tcal_f(t),}
\end{equation}
with $\eta_{\rm store}=25\Ocal/[16(kR)^2]$ for the aligned plus-quadrupole wave-zone model. The source is initialized by a local equal-charge encoder, and an exact controller-empty handoff permits the complete coherent branch information to be compressed into one virtual bosonic difference mode. This supplies the physical channel input required by the subsequent bosonic analysis rather than assuming it at the outset. Precursor radiation emitted during the encoder remains part of the complete waveform and is not reset at handoff.

The memory quantum excess is $\Delta_{\rm mem}=\tau_c-m_c$; an accessible readout further requires $\Delta_{\rm acc}=\tau_r\Delta_{\rm mem}-m_r>0$. In vacuum, the exact binary-coherent negativity shows $\mathcal N_{\max}\sim\tau_c$ for weak links. For two ordinary passive nonrelativistic matter interfaces, the explicit benchmark produces an end-to-end coherent scale near $10^{-42}$, revealing that source gravitational branching is as important as receiver capture.

The resulting framework is best viewed as a normalization benchmark for future gravitational quantum transducers. Any proposed improvement can be asked to identify which factor it changes, what new noise it introduces, and whether the improvement survives all the way to an accessible quantum register.

\appendix

\section{Branch-distance conservation at the encoder handoff}
\label{app:handoff}

For the damped two-mode encoder, the source amplitude is
\begin{equation}
\alpha_s(t)=-is\zeta\frac{g}{\Omega}\ee^{-\kappa_A t/4}\sin(\Omega t).
\end{equation}
The branch separation of the source mode is $2\alpha_+(t)$. Each Markov output port $j$ carries branch-difference amplitude $2\sqrt{\kappa_j}\alpha_+(t)$. The branch distance emitted into port $j$ through the controller-empty time is
\begin{equation}
\boxed{
N_{\Delta,j}^{\rm enc}
=4\frac{\kappa_j}{\kappa_A}|\zeta|^2
\left(1-\ee^{-\kappa_A T_*/2}\right).}
\end{equation}
The source branch distance remaining at handoff is
\begin{equation}
N_{\Delta,A}(T_*)=4|\zeta|^2\ee^{-\kappa_A T_*/2}.
\end{equation}
Summing all source output ports gives
\begin{equation}
N_{\Delta,A}(T_*)+
\sum_jN_{\Delta,j}^{\rm enc}=4|\zeta|^2,
\end{equation}
which is Eq.~\eqref{eq:distance-conservation}. In the fast-encoder limit, the emitted precursor fraction is
\begin{equation}
\epsilon_{\rm pre}
=1-\ee^{-\kappa_A T_*/2}
\simeq\frac{\pi\kappa_A}{4g}.
\end{equation}

\section{Exact pure-loss partial-transpose block}
\label{app:pureloss}

For Eq.~\eqref{eq:hybrid}, define normalized receiver even and odd cat states
\begin{equation}
|C_+\rangle=\frac{|B\rangle+|-B\rangle}{\sqrt{2(1+s_B)}},
\qquad
|C_-\rangle=\frac{|B\rangle-|-B\rangle}{\sqrt{2(1-s_B)}}.
\end{equation}
The partially transposed state decomposes into two $2\times2$ blocks. The block containing the negative eigenvalue is
\begin{equation}
M_-=
\frac14
\begin{pmatrix}
(1+s_B)(1-s_E) & \sqrt{1-s_B^2}(1+s_E)\\
\sqrt{1-s_B^2}(1+s_E) & (1-s_B)(1-s_E)
\end{pmatrix}.
\end{equation}
Its determinant is
\begin{equation}
\det M_-=-\frac{s_E(1-s_B^2)}{4}<0
\end{equation}
for every $A>0$ and $\eta>0$. The magnitude of its negative eigenvalue gives Eq.~\eqref{eq:exactN}.

\section{Passive quadrupole sum-rule bound}
\label{app:sumrule}

Let a stationary nonrelativistic matter system have Hamiltonian eigenstates $|m\rangle$ with energies $E_m$ and a passive diagonal state $\rho=\sum_m p_m|m\rangle\langle m|$, so $p_m\ge p_n$ whenever $E_m<E_n$. For any Hermitian operator $F$,
\begin{equation}
\frac12\Tr\rho[F,[H,F]]
=\sum_{m<n}(p_m-p_n)(E_n-E_m)|F_{mn}|^2.
\end{equation}
For the five orthonormal STF mass-quadrupole components $Q_A$, the coordinate-space commutator identity gives
\begin{equation}
\boxed{
\sum_A\frac12\Tr\rho[Q_A,[H,Q_A]]
=\frac{10}{3}\hbar^2\langle I\rangle_\rho,}
\end{equation}
where $I=\sum_a m_a r_a^2$. Passivity makes every positive-frequency term on the spectral side nonnegative. For a narrow selected band near $\omega$,
\begin{equation}
\sum_{\rm band,A}(p_m-p_n)|Q^A_{mn}|^2
\lesssim\frac{10}{3}\frac{\hbar\langle I\rangle}{\omega}.
\end{equation}
Using the quadrupole graviton transition rate
\begin{equation}
\gamma_{nm}^{(g)}
=\frac{2G\omega_{nm}^5}{5\hbar c^5}
Q_{ij}^{nm}Q_{ij}^{mn},
\end{equation}
one obtains the passive net gravitational linewidth ceiling
\begin{equation}
\kappa_{g,\rm net}
\lesssim\frac{4G}{3c^5}\langle I\rangle\omega^4.
\end{equation}
Writing $\langle I\rangle=ML^2$, $\Ccal=2GM/(c^2L)$, and $\beta=\omega L/c$ gives
\begin{equation}
\boxed{
\frac{\kappa_{g,\rm net}}{\omega}
\lesssim\frac23\Ccal\beta^3,}
\end{equation}
which is Eq.~\eqref{eq:sumrule}. The positivity step is exactly where the passive-state assumption enters; active or inverted states are not constrained by this positive-term argument.

\section{Passive source broadening}
\label{app:broadening}

For a passive exponential source with $r=\kappa_A/\kappa_B$, fixed intrinsic source gravitational rate $\kappa_{g,A}$, and fixed receiver parameters, the optimized end-to-end transfer is
\begin{equation}
\tau_{A\to B}^{\max}(r)
=4\frac{\kappa_{g,A}\kappa_\Delta}{\kappa_B^2}
r^{2r/(1-r)}.
\end{equation}
Since
\begin{equation}
\frac{\dd}{\dd r}\ln r^{2r/(1-r)}
=\frac{2(\ln r+1-r)}{(1-r)^2}<0,
\end{equation}
adding passive nongravitational source damping strictly reduces the source-resolved optimum. Receiver-local temporal matching and source branching must therefore not be optimized independently.

\bibliographystyle{unsrtnat}
\bibliography{references}

\end{document}
